\section{Conclusion}
\label{sec:conclusion}

This paper presents DeltaCycle, a capacity-only monitoring loop built on a recurrent linear-attention backbone, 
which transforms battery health monitoring from cloud-side accuracy benchmarking into a decision-grade capability 
that runs directly on the device. The backbone is built on a gated DeltaNet-2 layer with a fixed-size matrix state, 
which occupies only 8 KB per layer in the deployed configuration. 
This bounded state footprint is precisely what enables the same model weights to be executed directly on a 
BMS microcontroller. Degradation is captured at three time scales by patch branches (2/4/8) reading local, regional, 
and global views. A physics-consistent rate head anchors predictions to the last observed capacity  with a 
non-negative resistance coefficient that only accelerates decay.  A pinball-trained quantile head outputs 
P2.5/P50/P97.5 in a single forward pass, which conformal calibration on a held-out cell turns into a coverage-verified band.
We evaluate the proposed framework under the per-SP protocol of PatchFormer and RUL-Mamba on six public datasets 
spanning LCO, NCA, NCM, and LFP chemistries, with architecture, physics, uncertainty, 
and deployment studies validated under the same protocol discipline. The performance 
of DeltaCycle is validated across multiple publicly available datasets, 
leading to the following conclusions:

\begin{enumerate}
\item \textbf{First-tier trajectory accuracy across six degradation
regimes.} On mean per-SP MAE the model ranks top-3 on all six
datasets---first on GOTION, PANASONIC and TJU (tied), second on CALCE,
third on NASA and MIT---while clearly beating every generic time-series
transformer. On CALCE it ranks second to PatchFormer on
every reported metric (MAE $0.0065$--$0.0085$ vs.\ $0.0057$--$0.0081$), on
MIT it reaches $R^2 \ge 0.9995$ at every SP on both held-out cells, and on
NASA it exceeds the in-house PatchFormer in $R^2$ at SP70 and SP90
($0.9822$ vs.\ $0.9808$; $0.9804$ vs.\ $0.9623$) while holding a mean AE of
$0.6$ cycles against $2.5$. Non-monotonic field behaviour is tracked rather
than smoothed away: the model follows a $0.024$\,Ah local capacity recovery
on PANASONIC without overshoot, and the coarse branch encodes the
degradation stage itself ($93.2\%$ three-stage linear separability).

\item \textbf{Physics pays for itself where physics should matter.} A
systematic injection-route study first shows that the conventional routes
fail under the standard protocol which is
why the mechanism that does work moves the constraint into absolute capacity
space and onto an independently measured quantity. Trained that way, the
rate head lowers CALCE MAE by $53\%$ against a control free head sharing the
same target space. On the
unseen final $10\%$ of life it reaches $R^2 = 0.7745\pm0.0001$ where the
main-model readout reaches $0.374$ , more than doubling tail fidelity. Under a corrupted capacity
channel it wins every robustness setting on MAE. The constraint is interpretable as well as effective: the
resistance coefficient is non-negative by construction, so measured
resistance can only accelerate the predicted decay and the head suppresses non-physical regeneration artifacts.

\item \textbf{Uncertainty that a maintenance rule can consume.} Raw pinball
quantiles cover only $89.1\%$ of observations at a nominal $95\%$; conformal
calibration on a held-out cell, using two scalar parameters and no
ensembling, lifts coverage to $93.5\%\pm0.9\%$. The residual gap to the nominal level reflects
finite-sample calibration and cross-cell drift, and is reported rather than
hidden. 

\item \textbf{Deployment verified at the level firmware reviews require.}
The full multi-scale model in freestanding C99 matches the PyTorch forward
to $1.07\times10^{-6}$ relative error, and the Cortex-M3 binary matches the
host binary to $5.96\times10^{-8}$---about half a float32 ULP. INT8
quantization takes the weights from $1.85$\,MB to $504$\,KB with no
measurable accuracy loss ($R^2$ $0.9955$ unchanged, AE $2.16 \to 2.17$
cycles), and the minimal single-branch variant fits a $128$\,KB part at
$122$\,KB. The recurrent state stays at $48$\,KB regardless of window
length. The stack budget is $3.1$\,KB and fixed at compile time.
\end{enumerate}

DeltaCycle shows that trajectory accuracy, physical consistency, calibrated uncertainty,
and on-device feasibility are not mutually exclusive and can be achieved simultaneously.
Looking ahead, our future work will focus on two fronts. Architecturally, we will
refine the cross-scale exchange to make the rate head adaptive to cells whose impedance
channel is uninformative. On the deployment side, we will push compression beyond INT8---mixed
precision and pruning---so that the multi-scale model fits the smaller flash parts
used in volume BMS designs.
